% TOP_PROBLEM: {"id": 3103, "title": "Reconstruction conjecture", "category_id": 3, "category": {"id": 3, "name": "graph_theory", "display_name": "Graph Theory", "description": "Problems involving graphs, networks, and their properties.", "slug": "graph-theory", "order_index": 3, "created_at": "2026-07-31T15:26:25.671Z"}, "status": "open", "problem_number": "OPG-658", "set_id": 12}
% FIRST_POSTED: 2026-10-01
% ENTRY_KIND: statement_only
% UnsolvedMath ID 3103; source catalogue code OPG-658
% !TeX program = lualatex
% Definition and sources only; this file is not an LLM solution attempt.
\documentclass[11pt,a4paper]{article}
\usepackage[margin=25mm]{geometry}
\usepackage{fontspec}
\setmainfont{lmroman10-regular.otf}[BoldFont=lmroman10-bold.otf,
 ItalicFont=lmroman10-italic.otf,BoldItalicFont=lmroman10-bolditalic.otf]
\usepackage{amsmath,amssymb,mathtools}
\usepackage{unicode-math}
\setmathfont{latinmodern-math.otf}
\AtBeginDocument{\let\setminus\smallsetminus}
\usepackage{xurl,hyperref}
\hypersetup{colorlinks=true,urlcolor=blue}
\setcounter{secnumdepth}{0}
\setlength{\parindent}{0pt}
\setlength{\parskip}{5pt}
\setlength{\emergencystretch}{3em}
\begin{document}
\section*{Reconstruction conjecture}\label{3103}
{\small UnsolvedMath ID 3103 · OPG-658 · Graph Theory · OpenGarden}
\subsection{Definitions and mathematical statement}
Let $G=(V,E)$ be a finite simple undirected graph. For $v\in V$, the
vertex-deleted graph $G-v$ is the induced graph on $V\setminus\{v\}$:
the vertex and all its incident edges are removed. Write $[H]$ for the
isomorphism class of an unlabelled graph $H$, where graph isomorphism is
a vertex bijection preserving adjacency in both directions.

The vertex deck of $G$ is the multiset
\[
 \mathcal D(G)=\{\!\{[G-v]:v\in V\}\!\}.
\]
There is one card for each vertex, and repeated isomorphism classes are
retained with their multiplicities. Neither the deleted vertex's identity
nor a consistent labelling across the cards is supplied.

The reconstruction conjecture states that for any finite simple graphs
$G,H$ with at least three vertices,
\[
 \mathcal D(G)=\mathcal D(H)\quad\Longrightarrow\quad G\cong H.
\]
Equality of decks includes equality of multiplicities, hence their
cardinalities and the orders of the original graphs. Both connected and
disconnected graphs, including isolated vertices, are allowed.

The claim is uniqueness up to isomorphism, not recovery of an original
vertex labelling. A counterexample would be a pair of nonisomorphic
simple graphs on the same number of at least three vertices with identical
decks. Matching a selection of cards, or merely matching counts of small
subgraphs, is weaker than deck equality. Directed graphs, parallel edges
and loops are outside this formulation, since altering the graph class
changes the reconstruction question.
\subsection{Short English statement}
Is every finite simple graph with at least three vertices uniquely determined up to isomorphism by its multiset of vertex-deleted graphs?
\subsection{Sources}
\begin{itemize}
\item[S1] ULAM.AI and the UnsolvedMath Contributors, supplied problems.json, record 3103 (OPG-658), OpenGarden collection. Catalogue identity and source transcription; CC BY 4.0.
\url{https://huggingface.co/datasets/ulamai/UnsolvedMath}.
\item[S2] Open Problem Garden, Reconstruction conjecture. Original problem and discussion; the mathematical formulation below is expanded from the supplied source record.
\url{http://www.openproblemgarden.org/op/reconstruction_conjecture}.
\end{itemize}
\end{document}
